\section{Campagnes de test --- environnement géographique}
\subsection{Buildings}
\begin{tests}
\T{BLD-01}{P2}{Dessiner un bâtiment (clics, clic droit pour fermer\man{29}) ; renseigner \ui{Text}, \ui{Building height}, \ui{Base elevation}.}{Le dernier segment se ferme automatiquement. Libellé et hauteur affichés sur le plan.}
\T{BLD-02}{P1}{Avec l'origine de \code{BASE-0}, importer \code{buildings\_ok.kml} (\ui{Import building polygons from a KML file}\man{29,31}).}{\textbf{[M]} nom = \code{hauteur;libellé}\man{31}. Attendu : \code{10;B01} $\to$ 10~m « B01 » ; \code{12.5;B02} $\to$ 12,5~m ; \code{15;L-shape} (polygone concave en L) ; \code{4;2} $\to$ 4~m « 2 » (exemple du manuel\man{32}). 4 bâtiments, positions à $\pm0{,}5$~m des coordonnées du KML (origine du manuel).}
\T{BLD-03}{P1}{Importer \code{buildings\_edge.kml} (19 éléments). Pour chacun noter : importé / ignoré / corrigé, hauteur et libellé lus. Tableau de lecture : annexe~\ref{sec:kmledge}.}{Aucun plantage ni blocage. \textbf{[O]} sauf : les éléments sans polygone (\code{Point}, \code{LineString}) ne doivent \emph{pas} devenir des bâtiments ; \code{12,5} (virgule) ne doit pas devenir \code{125} ; hauteur $\le0$ ou $10^4$~m ne doit pas fausser un calcul sans message.}
\T{BLD-04}{P1}{Importer successivement \code{malformed\_truncated}, \code{empty}, \code{utf8\_bom}, \code{crlf}, \code{utf16}, puis renommer un fichier texte en \code{.kml}.}{Corrompu/vide : message et projet intact. BOM/CRLF/UTF-16 : \textbf{[O]}. Un import qui échoue n'altère pas les bâtiments existants.}
\T{BLD-05}{P1}{Réimporter \code{buildings\_ok.kml} deux fois de suite.}{\textbf{[O]} doublons ou fusion ? Noter le nombre final (4 ou 8).}
\T{BLD-06}{P2}{Importer \code{buildings\_stress\_2000.kml}.}{Temps d'import, fluidité du zoom, rapport et 3D (\code{ROB-05}). Pas de plantage.}
\T{BLD-07}{P1}{OSM\man{30,31} : charger la zone visible. Boîte de hauteurs : 20~m par défaut, 3~m par étage, altitude de base 0\man{31}.}{\textbf{[M]} bâtiment à \code{building:levels=4} sans hauteur $\to$ $4\times3=12$~m. Sans niveaux ni hauteur $\to$ 20~m. Charger la même zone une 2\textsuperscript{e} fois : \textbf{aucun doublon} et hauteurs éditées conservées\man{31}.}
\T{BLD-08}{P2}{Zone OSM très large (zoom 12), puis hors ligne.}{Avertissement conseillant de ne pas charger plus que nécessaire\man{31}. Hors ligne : message, pas de gel.}
\T{BLD-09}{P2}{Copier, supprimer (un/tous), ajouter/supprimer un n\oe ud, \ui{Square building corners}, envoyer à l'arrière, flèches de déplacement d'un pixel, annuler/rétablir\man{29}.}{Angles à $90^\circ\pm1^\circ$ après mise à l'équerre. Annuler/rétablir fonctionnent sur 20 opérations et après « supprimer tous ».}
\T{BLD-10}{P1}{\ui{Fill in all buildings base elevation}\man{29} : bâtiment au-dessus de points de terrain 100/102/104/106 sur ses 4 coins ; choisir min, max, moyenne.}{\textbf{[F]} 100 ; 106 ; 103~m.}
\T{BLD-11}{P1}{\textbf{Blindage complet}\man{30} : importer \code{shield\_scenario.kml} (mur 50~m de haut de $x=30$ à 50, $y\in[-30,30]$). Source \code{TS07} à (0,0), hauteur d'antenne portée à 20~m ; point de contrôle à (100,0) hauteur 1,5~m. Modes : \ui{Do not use}, \ui{Complete shielding}.}{\textbf{[M]} sans bâtiments : $E>0$. Blindage complet : $E=0$\man{30}.}
\T{BLD-12}{P1}{Même scénario, mode blindage, hauteurs de point 60, 100, 119, 121~m.}{\textbf{[F géométrie]} le rayon qui monte de $(0;20)$ vers $(100;h)$ passe au-dessus de la face avant du mur ($x=30$, haut à 50~m) si $20+0{,}3(h-20)\ge50$, soit $h\ge120$~m. Attendu : $E=0$ à 60, 100 et 119~m ; $E>0$ à 121~m.}
\T{BLD-13}{P1}{\textbf{Atténuation interne}\man{30} (6~dB sous 500~MHz, 3~dB au-dessus\man{52}). Point de type \code{Building} dans \code{INSIDE}. Source à 900~MHz puis à 150~MHz. Comparer à \ui{Do not use buildings}.}{\textbf{[R]} rapport $E_{\mathrm{in}}/E_{\mathrm{out}}=10^{-A/20}$ : 0,708 à 900~MHz (3~dB), 0,501 à 150~MHz (6~dB). Des rapports 0,501 (3~dB) et 0,251 (6~dB) indiqueraient que l'affaiblissement est appliqué au champ comme s'il s'agissait d'une puissance. \textbf{[O]} source à exactement 500~MHz.}
\T{BLD-14}{P2}{Décocher \ui{Take into account custom attenuation inside building} pour un seul bâtiment\man{29} ; décocher \ui{Show on the horizontal cut}.}{Ce bâtiment n'applique plus l'atténuation interne ; il disparaît des coupes horizontales mais reste pour le calcul.}
\T{BLD-15}{P1}{Points \code{Ground} et \code{Building}\man{51} avec une altitude de base de 100~m, hauteur de point 1,1~m.}{\textbf{[M]} \code{Building} : hauteur mesurée depuis la base du bâtiment ; \code{Ground} : depuis le terrain.}
\T{BLD-16}{P2}{Polygones qui se chevauchent, imbriqués, identiques ; polygone à un côté de 0,5~m.}{Pas de plantage. \textbf{[O]} règle d'application du blindage.}
\T{BLD-17}{P1}{\textbf{Outil « Delete objects outside the screen »}\man{58} (2 variantes). Compter les objets avant, puis après.}{Objets hors écran supprimés, objets dans l'écran conservés (variante 1 : terrain compris ; variante 2 : sans terrain). \textbf{[O]} bâtiment à cheval sur le bord. Y a-t-il une confirmation et un \emph{Undo} ?}
\end{tests}

\subsection{Terrain}
\begin{tests}
\T{TER-01}{P1}{Placer 3 points : (0,0)=100~m, (100,0)=110~m, (0,100)=100~m. Afficher la triangulation\man{34}, lire l'altitude au pointeur (barre d'état).}{\textbf{[F]} plan $z=100+0{,}1x$ : 105,0~m en (50,0) ; 102,5~m en (25,25) ; 100,0~m en (0,50). Interpolation linéaire à $\pm0{,}05$~m.}
\T{TER-02}{P2}{Ajouter une ligne de niveau (clic droit pour finir, double clic droit pour revenir à la sélection\man{34}), puis points, déplacer, supprimer.}{La triangulation se refait toute seule en quittant le panneau\man{34}.}
\T{TER-03}{P1}{Charger un MNT\man{33} (en ligne).}{\textbf{[M]} tout le terrain existant est remplacé\man{34} (message ?). Les points MNT sont ensuite modifiables comme des points manuels.}
\T{TER-04}{P1}{\textbf{Invariance à l'altitude} : calculer \code{PNT-03} ; ajouter +100~m à tous les points de terrain et altitudes de base ; recalculer.}{\textbf{[R]} résultats identiques (hauteurs relatives au sol\man{15,51}).}
\T{TER-05}{P1}{Tour de 30~m sur une colline (plan de \code{TER-01}) ; coupe verticale.}{Le profil du terrain est tracé\man{38} ; l'antenne se trouve à 30~m \emph{au-dessus du terrain au pied de la tour}\man{15}.}
\T{TER-06}{P2}{Sans aucun point de terrain.}{\textbf{[M]} terrain plat à 0 ; on ne peut pas le désactiver\man{32}.}
\T{TER-07}{P1}{Cas limites : 0, 1, 2 points ; 3 points alignés ; points en double ; altitude 9000 ; $-500$ ; texte ; contour à un seul n\oe ud ; contours qui se croisent ; point hors grille.}{Pas de plantage ni de triangulation fantôme. Messages explicites quand la triangulation est impossible.}
\T{TER-08}{P2}{Copier, supprimer (un/tous), annuler/rétablir\man{33,34}.}{Comportement cohérent avec \code{BLD-09}.}
\end{tests}

\subsection{Lines and Labels --- routes, tours, repères}
\begin{tests}
\T{LBL-01}{P2}{Dessiner route et barrière ; largeur de route et épaisseur ; charger depuis OSM\man{35,36}.}{Objets créés, paramètres appliqués. OSM hors ligne : message.}
\T{LBL-02}{P3}{Étiquette : texte, police, taille, icône, rotation\man{36,38}.}{Tous les paramètres modifient l'étiquette.}
\T{LBL-03}{P1}{Créer l'étiquette \code{Tower/15/15/0.5/1.2} (exemple du manuel\man{37,38}) puis des vertical cuts passant par la tour.}{\textbf{[M]} format \code{Tower/H1/H2/L1/L2}. La tour apparaît en coupe verticale\man{38} et en 3D. $H2=0$ $\to$ prisme rectangulaire\man{37}. Si la tour est dans le contour d'un bâtiment, elle est sur son toit.}
\T{LBL-04}{P2}{Variantes de tour (bloc \code{L-1}) : \code{Tower/15/0/0.5/0.5}, 3 paramètres, \code{tower/...} en minuscules, valeurs négatives ou nulles, virgule décimale, espaces, deux tours superposées, tour loin des sources.}{\textbf{[O]} règle de lecture. Une étiquette mal formée devient un simple texte, jamais un plantage de la coupe.}
\T{LBL-05}{P3}{N\oe uds de ligne : ajouter, supprimer, couper (\ui{Break knot})\man{35}. Flèches de déplacement d'un pixel.}{Fonctionnent ; annulation possible.}
\end{tests}

\subsection{Calibrate Map Image et fonds de carte}
\begin{tests}
\T{CAL-01}{P2}{Charger \code{calib\_grid} en \code{.png}, \code{.jpg}, \code{.bmp}, \code{.gif}, \code{.tif}\man{22}.}{Les 5 formats se chargent, couleurs fidèles.}
\T{CAL-02}{P1}{Charger les images pièges : \code{edge\_1x1}, \code{edge\_9000x9000}, \code{edge\_transparent}, \code{edge\_grayscale}, \code{edge\_cmyk.jpg}, \code{edge\_corrupt}, \code{edge\_empty}, \code{edge\_text\_renamed}.}{Aucun plantage. Les trois dernières : message d'erreur et ancienne image conservée. \code{9000$\times$9000} : \textbf{[O]} mémoire et temps.}
\T{CAL-03}{P1}{Méthode « point de référence + ligne »\man{22}. Sur \code{calib\_grid.png} : ligne rouge \code{REF A} de $(100,100)$ à $(600,100)$~px, longueur \textbf{100~m} ; point de référence sur le même endroit de la carte OSM. Choisir \ui{Calibrated map image} puis mesurer avec la règle.}{\textbf{[F]} 0,2~m/px. \code{CHECK B} (300~px) $=60$~m ; \code{CHECK C} (500~px, verticale) $=100$~m, à $\pm1\,\%$ (limité par le pointage).}
\T{CAL-04}{P1}{Méthode « deux lignes »\man{23}. Image \code{calib\_grid\_rot30.png} (tournée de $30^\circ$) ; ligne de référence image = \code{REF A} ; ligne carte = segment de 100~m orienté à l'azimut $90^\circ$ (est).}{\textbf{[F]} après calibrage, \code{REF A} est est-ouest et \code{CHECK C} est nord-sud ($0^\circ$/$180^\circ\pm1^\circ$) : la rotation est corrigée.}
\T{CAL-05}{P2}{Recalibrer. Sans puis avec \ui{Refresh the pyramid of tiles}\man{23}. \ui{Delete image}.}{Sans rafraîchissement l'ancien calibrage peut rester affiché ; avec, le nouveau apparaît\man{23}. \textbf{[O]}.}
\T{CAL-06}{P1}{Longueur de référence 0, $-100$, texte, $10^9$ ; ligne de 0~px ; ligne de 1~px.}{Rejet clair. Jamais d'échelle infinie ni \code{NaN} sur la carte.}
\T{CAL-07}{P2}{Enregistrer, fermer, rouvrir un projet avec image calibrée.}{Calibrage identique ; taille du \code{.xemf} cohérente avec l'image.}
\T{CAL-08}{P1}{\textbf{Fonds de carte personnalisés} : lancer \code{python3 tools/tile\_server.py}. Dans \ui{Settings} $\to$ \ui{Base map settings}, ajouter \code{LOCAL\_XYZ} = \code{http://127.0.0.1:8000/t/}\allowbreak\code{[Z]/[X]/[Y].png} et \code{LOCAL\_QUERY} = \code{http://127.0.0.1:8000/g?x=[X]}\allowbreak\code{\&y=[Y]\&zoom=[17-Z]}\man{60,61}. Appliquer, choisir la carte.}{\textbf{[F]} les tuiles affichent leur \code{z}, \code{x}, \code{y}. Le script journalise chaque requête. \code{[17-Z]} doit être évalué en expression. À l'échelle 18 la tuile contenant l'origine du manuel est $x=137770$, $y=93797$ (19 : 275540, 187595). Une inversion $X/Y$ ou un décalage d'un niveau de zoom se voit immédiatement.}
\T{CAL-09}{P2}{Fond de carte invalide : URL sans \code{[Z]}, serveur arrêté, HTTP 404/500, délai dépassé, mauvais proxy. \ui{Reset to default} puis \ui{Apply base maps settings}\man{60}.}{Messages ou cases vides, pas de gel de l'interface. Réinitialisation fidèle à la liste d'origine.}
\T{CAL-10}{P3}{Chaque carte de la liste\man{60} (OSM, Carto, OSM Relief, Google, Esri, HERE, Mapquest, OpenTopoMap\ldots) à 2 ou 3 zooms.}{Les tuiles s'affichent. \textbf{Respectez les conditions d'usage des serveurs} : quelques tuiles seulement, pas de balayage massif. Vérifier la présence de Bing, citée p.4 (\code{DOC-13}).}
\T{CAL-11}{P2}{Dossier de cache : tuiles créées ; passer hors ligne et revisiter la zone\man{60,61}.}{Les tuiles déjà vues restent affichées sans réseau.}
\T{CAL-12}{P3}{Menu zoom et roulette (clic sur la roulette pour déplacer\man{9}). Changer de fond de carte et de zoom entre deux calculs.}{\textbf{[R]} aucun résultat chiffré ne change quand on change de fond ou de zoom.}
\end{tests}
